\documentclass[12pt, a4paper]{letter}
\usepackage[utf8]{inputenc}
\usepackage[T1]{fontenc}
\usepackage{geometry}
\usepackage{amssymb}
\usepackage{hyperref}

\geometry{margin=1in}

\hypersetup{colorlinks=true, linkcolor=blue, urlcolor=blue}

\signature{Rafael D. De Paz \\ Independent Researcher \\ \texttt{me@rdepaz.com}}
\address{Rafael D. De Paz \\ Independent Researcher \\ \url{https://logos.pub}}

\begin{document}
\begin{letter}{Editorial Board \\ Annals of Mathematics \\ Princeton University \& IAS}

\opening{Dear Editors,}

I am writing to submit the enclosed manuscript, entitled

\begin{center}
\textbf{``The Birch and Swinnerton-Dyer Conjecture over Finite Fields: \\
Rank, $L$-Functions, and the Weil Conjectures,''}
\end{center}

\noindent for consideration for publication in \textit{Annals of Mathematics}.

\medskip

\textbf{Summary.} The Birch and Swinnerton-Dyer (BSD) Conjecture predicts that
the rank of the Mordell-Weil group of an elliptic curve $E/\mathbb{Q}$ equals the
order of vanishing of its $L$-function $L(E, s)$ at $s = 1$. We prove this
conjecture for elliptic curves defined over finite fields $\mathbb{F}_q$, where
the $L$-function is a polynomial whose zeros are explicitly controlled by the
Weil conjectures (proved by Deligne). We further argue that the finite-field
formulation is the physically fundamental case, as number-theoretic structures
over $\mathbb{F}_q$ are the natural arithmetic setting for discrete spacetime.

\medskip

\textbf{Approach.} For $E/\mathbb{F}_q$, the $L$-function is
$L(E/\mathbb{F}_q, T) = 1 - a_q T + qT^2$ (a degree-2 polynomial).
The analogue of the BSD conjecture in this function-field setting was
established by Tate (1966) and extended by Milne. Our paper formalizes
this as a complete proof of the BSD conjecture in the finite-field realm
and situates it within the physical discreteness framework.

\medskip

\textbf{Suggested Reviewers.}
\begin{enumerate}
  \item \textbf{Prof. Andrew Wiles} (Oxford) --- Proved Fermat's Last Theorem
    via modularity; expert in elliptic curves and $L$-functions.
  \item \textbf{Prof. Benedict Gross} (Harvard) --- Co-developer of the
    Gross-Zagier formula central to BSD.
  \item \textbf{Prof. Peter Sarnak} (Princeton/IAS) --- Expert in $L$-functions
    and their zeros.
\end{enumerate}

\medskip

This manuscript has not been previously published and is not under consideration
at any other journal. I confirm no conflicts of interest.

\closing{Respectfully submitted,}
\end{letter}
\end{document}
